\section{Analyse fonctionnelle}
\subsection{Vision globale du système}
% TODO: Insérer et commenter le Schéma Conceptuel de l'infrastructure globale.
% \begin{figure}[htbp]
%   \centering
%   \include{diagrams/schema_conceptuel}
%   \caption{Architecture globale des flux entre le Client, le Résolveur et la Hiérarchie DNS.}
% \end{figure}
% - Expliquer le trajet d'une requête de bout en bout (Client -> Résolveur -> Racine -> TLD -> Autorité).

\subsection{Épic 1 : L'Outil Client (CLI)}
% TODO: Insérer le Use Case Diagram du Client.
% - Lister et détailler les US-1.1 à US-1.6 (envoi de requêtes, parsing d'arguments, affichage formaté).

\subsection{Épic 2 : Le Serveur Récursif (Le Chef d'orchestre)}
% TODO: Insérer le Use Case Diagram du Résolveur Récursif.
% - Lister et détailler les US-2.1 à US-2.7 (moteur itératif, cache, fallback TCP).
% - Explain administration feature: interception de la demande 'flush.cache.local' (TXT).
% % [À MENTIONNER]: Règle métier stricte : rejet systématique des requêtes reçues avec RD=0 via le code RCODE 5 (REFUSED).

\subsection{Épic 3 : Le Serveur d'Autorité (La Base de données)}
% TODO: Insérer le Use Case Diagram du NameServer.
% - Lister et détailler les US-3.1 à US-3.6 (chargement de zone, réponses autoritaires, hot-reload).

\subsection{Épics 4 et 5 : Tolérance aux pannes et Qualité}
% TODO: Présenter les fonctionnalités avancées.
% - Épic 4 (US-4.1 à US-4.5) : Mécanismes de réplication Master/Slave (AXFR et NOTIFY).
% - Épic 5 (US-5.1) : Exigences de qualité et couverture par tests unitaires.